\section*{Acknowledgements}

Author contributions: Aditrivaishnavi Balaji conceived and led the project,
formulated the research questions, designed the study, developed and validated
the software and injection--recovery pipeline, curated the data and verified
its provenance, designed and executed the experiments and statistical
analyses, interpreted the results, produced the figures, tables, and
reproducibility materials, and drafted and revised the manuscript. Parker
T. Johnson supervised and mentored the project, advised on the astrophysical
context, methodology, and scientific interpretation, and critically reviewed
the manuscript. Both authors reviewed and approved the final manuscript.

The Legacy Surveys consist of three individual and complementary projects:
the Dark Energy Camera Legacy Survey (DECaLS; Proposal ID 2014B-0404; PIs:
David Schlegel and Arjun Dey), the Beijing--Arizona Sky Survey (BASS; NOAO
Proposal ID 2015A-0801; PIs: Zhou Xu and Xiaohui Fan), and the Mayall
\(z\)-band Legacy Survey (MzLS; Proposal ID 2016A-0453; PI: Arjun Dey).
DECaLS, BASS, and MzLS together include data obtained, respectively, at the
Blanco telescope, Cerro Tololo Inter-American Observatory, NSF's NOIRLab; the
Bok telescope, Steward Observatory, University of Arizona; and the Mayall
telescope, Kitt Peak National Observatory, NOIRLab. Pipeline processing and
analyses of the data were supported by NOIRLab and the Lawrence Berkeley
National Laboratory (LBNL). The Legacy Surveys project is honoured to be
permitted to conduct astronomical research on Iolkam Du'ag (Kitt Peak), a
mountain with particular significance to the Tohono O'odham Nation.

NOIRLab is operated by the Association of Universities for Research in
Astronomy under a cooperative agreement with the National Science Foundation.
LBNL is managed by the Regents of the University of California under contract
to the U.S. Department of Energy.

This project used data obtained with the Dark Energy Camera (DECam), which was
constructed by the Dark Energy Survey (DES) collaboration. Funding for the DES
Projects has been provided by the U.S. Department of Energy, the U.S. National
Science Foundation, the Ministry of Science and Education of Spain, the
Science and Technology Facilities Council of the United Kingdom, the Higher
Education Funding Council for England, the National Center for Supercomputing
Applications at the University of Illinois at Urbana--Champaign, the Kavli
Institute of Cosmological Physics at the University of Chicago, Center for
Cosmology and Astro-Particle Physics at the Ohio State University, the
Mitchell Institute for Fundamental Physics and Astronomy at Texas A\&M
University, Financiadora de Estudos e Projetos, Fundacao Carlos Chagas Filho
de Amparo a Pesquisa do Estado do Rio de Janeiro, Conselho Nacional de
Desenvolvimento Cientifico e Tecnologico and the Ministerio da Ciencia,
Tecnologia e Inovacao, the Deutsche Forschungsgemeinschaft and the
Collaborating Institutions in the Dark Energy Survey. The Collaborating
Institutions are Argonne National Laboratory, the University of California at
Santa Cruz, the University of Cambridge, Centro de Investigaciones
Energeticas, Medioambientales y Tecnologicas--Madrid, the University of
Chicago, University College London, the DES--Brazil Consortium, the University
of Edinburgh, the Eidgenossische Technische Hochschule (ETH) Zurich, Fermi
National Accelerator Laboratory, the University of Illinois at
Urbana--Champaign, the Institut de Ciencies de l'Espai (IEEC/CSIC), the
Institut de Fisica d'Altes Energies, Lawrence Berkeley National Laboratory,
the Ludwig Maximilians Universitat Munchen and the associated Excellence
Cluster Universe, the University of Michigan, NSF's NOIRLab, the University
of Nottingham, the Ohio State University, the University of Pennsylvania, the
University of Portsmouth, SLAC National Accelerator Laboratory, Stanford
University, the University of Sussex, and Texas A\&M University.

BASS is a key project of the Telescope Access Program, funded by the National
Astronomical Observatories of China, the Chinese Academy of Sciences (Strategic
Priority Research Program ``The Emergence of Cosmological Structures'', grant
XDB09000000), and the Special Fund for Astronomy from the Ministry of Finance.
BASS is also supported by the External Cooperation Program of the Chinese
Academy of Sciences (grant 114A11KYSB20160057) and the Chinese National Natural
Science Foundation (grants 12120101003 and 11433005).

The Legacy Survey team makes use of data products from the Near-Earth Object
Wide-field Infrared Survey Explorer (NEOWISE), a project of the Jet Propulsion
Laboratory/California Institute of Technology funded by the National
Aeronautics and Space Administration. Legacy Surveys imaging of the DESI
footprint is supported by the Director, Office of Science, Office of High
Energy Physics of the U.S. Department of Energy under contract
DE-AC02-05CH1123; the National Energy Research Scientific Computing Center, a
DOE Office of Science User Facility under the same contract; and the U.S.
National Science Foundation, Division of Astronomical Sciences, under contract
AST-0950945 to NOAO.

This work used \textsc{PyTorch} and \textsc{torchvision} for neural-network
training and inference; \textsc{NumPy}, \textsc{SciPy}, \textsc{pandas}, and
\textsc{pyarrow} for numerical and tabular analysis; \textsc{scikit-learn} for
classification metrics and linear probes; \textsc{healpy} for HEALPix
grouping; \textsc{Astropy} for coordinates; \textsc{GalSim} for COSMOS source
construction; and \textsc{Matplotlib} for figures.

AI language models (Claude, Anthropic; ChatGPT, OpenAI) were used as coding
assistants during the development of analysis scripts and for manuscript
copy-editing. All AI-assisted code was reviewed, tested, and verified by the
authors, who are solely responsible for the scientific content of this paper.

\section*{Data availability}

DESI Legacy Imaging Surveys DR10 imaging and catalogues are publicly available
at \url{https://www.legacysurvey.org/dr10/}. The GalSim COSMOS 25.2 source
products are the archived dataset of \citet{Mandelbaum2012},
\url{https://doi.org/10.5281/zenodo.3242143}.
The configurations, analysis code, compact score and feature tables,
out-of-fold predictions, and provenance records used to inspect the reported
calculations are available at
\url{https://github.com/aditrivaishnavi/stronglens-realism-gap}. A public,
versioned archival release of the remaining redistributable project products
will be deposited no later than publication, and its DOI will be added to the
final article. It will include the exact primary EfficientNetV2-S checkpoint
(SHA-256
\texttt{5d1082839aac7a0b\allowbreak
fb89342f84a9a533\allowbreak
b22e1111c15cb471\allowbreak
1255ed1b524a2b8d}), the 12-checkpoint panel manifest and hashes, and the
compact panel score and result tables. If repository size limits require it,
the checkpoint weights will occupy separate DOI-linked records. Because of
their size, the generated image banks will be archived separately from the
compact analysis products.

Third-party DR10 and COSMOS pixels remain subject to their source archives and
licences rather than being redistributed as original data by this work.

\section*{Conflict of interest}

The authors declare no conflict of interest.

\section*{Funding}

The authors declare no funding.
